\scanpage{3}
\uline{Proposition} \unsure{1|!} Soient $S$ un schéma, $A$ un schéma abélien sur $S$, $X$ un schéma sur $S$, \deleted{\unsure{tel que pour toute|tel que pour tout} \unsure{fibre géométrique}} \inserted{$L$ un faisceau inversible sur} $A\times_S X$, \deleted{Considérons l'homomorphisme} trivial sur $e_A\times_S X$, et tel que \deleted{\unsure{pour que tout|pour tout} \unsure{$L_x$} \illegible{1}} pour tt $x\in X$, le faisceau $L_x$ sur $A_x=A_s\otimes_{k(s)}k(x)$ ($s$ le pt de $S$ sous $x$) soit alg. équivalent à zéro. Considérons l'application
% slanted
\margin{peut se généraliser au cas $A$ un schéma en groupes \inserted{\unsure{plat|plats}}, $L$ sur $A\times_S X$, satisfaisant $(\pi_A\times\mathrm{id}_X)^*(L)\simeq\mathrm{pr}_{13}^*(L)\otimes\mathrm{pr}_{23}^*(L)$ et $F$ et $G$ sur $A$ à supports propres sur $S$ et \deleted{\unsure{parfait}} \inserted{$F,G$} $\in D^{b}_{\mathrm{qcoh}}(A)$\ldots}
\[
\varphi_L\colon K^{\cdot}(A)\longrightarrow K^{\cdot}(X)
\]
définie par \deleted{l'application additive} le foncteur de catégories triangulées \deleted{$\unsure{\varphi}$} $\Phi\colon D_{\mathrm{parf}}(A)\unsure{\mathrel{-}|\longrightarrow}D_{\mathrm{parf}}(X)$
\[
\Phi_L(F)=\mathbb{R}\mathrm{pr}_{2*}\bigl(\mathbb{L}\mathrm{pr}_1^*(F)\otimes L\bigr).
\]
Alors $\varphi_L$ transforme produit de Pontryagin $*$ en produit ordinaire. Plus précisément, on a un isomorphisme de foncteurs de catégories triangulées
\[
\Phi_L(F\overset{\mathbb L}{*}G)\simeq\Phi_L(F)\overset{\mathbb L}{\otimes}\Phi_L(G)
\]
$(F,G\in\mathrm{Ob}\,D_{\mathrm{parf}}(A))$.
% slanted
\margin{on peut enfin généraliser en un énoncé concernant \deleted{\unsure{un|des} \unsure{couple|complexes}} de $C$-torseurs \diagram{\begin{tikzcd}[ampersand replacement=\&,column sep=tiny,row sep=tiny]P\times Q\arrow[r] \& R\\ C\times C\arrow[r] \& C\end{tikzcd}}, et \unsure{même} on peut laisser tomber $C$, et partir de faisceaux $L,L',L''$ sur $R\times X,P\times X,Q\times X$, et $\pi\colon P\times Q\unsure{\mathrel{-}|\to} R$, \deleted{\unsure{tels que}} tels que $\deleted{\unsure{R^{*}}}(\pi\unsure{\times|\mathbin{\&}|\alpha}\mathrm{id}_X)^*(L)\simeq\mathrm{pr}_{13}^*(L')\,\mathrm{pr}_{\unsure{22|23}}^*(L'')$}

Si le th. est démontré dans la situation donnée, il en résulte qu'il est vrai en tte situation \inserted{$(A,X',L')$} qui s'en déduit par image inverse, à l'aide d'un $S$-morphisme $f\colon X'\unsure{\mathrel{-}|\to}X$, car $\Phi_{L'}\simeq\mathbb{L}f^*(\Phi_L)$, et $\mathbb{L}f^*$ commute à $\overset{\mathbb L}{\otimes}$.

Soit $B$ le schéma abélien dual de $A$, et $L_0$ le faisceau de Weil sur $A\times_S B$. Alors l'hypothèse sur $L$ implique que $L$ est $\simeq$ l'image inverse de $L_0$ par un $S$-morphisme $X\longrightarrow B$. Donc on est ramené au cas où $X=B$, $L=L_0$. \deleted{\unsure{Or|On}} On a des isom. fonctoriels
\[
F\overset{\mathbb L}{\deleted{\otimes}{*}}G=\mathbb R\pi_{\unsure{*|\alpha}}\bigl(F\overset{\mathbb L}{\otimes}_S G\bigr)\qquad\text{et}\tag*{$(*)$}
\]
% archive sheet number: 2
